O sistema Mamdani apresentou decisões qualitativas coerentes com o conhecimento comum sobre conforto térmico~\cite{ross2010fuzzy}. A escolha dos operadores lógicos (mínimo, máximo, produto, soma) e das técnicas de defuzzificação (centroide, média do máximo, bissetriz) influenciou levemente o valor final da saída, mas a tendência qualitativa se manteve. O método centroide mostrou-se mais estável e sensível a todas as ativações das regras~\cite{ross2010fuzzy}.

A modelagem com funções de pertinência triangulares permitiu uma interpretação intuitiva das regras e dos resultados. O sistema é flexível e pode ser facilmente ajustado para diferentes ambientes ou critérios de conforto.

O sistema Takagi-Sugeno foi capaz de aproximar a função alvo com boa precisão, especialmente na versão de primeira ordem~\cite{takagi1985fuzzy}. O uso de funções de pertinência gaussianas proporcionou suavidade na transição entre as regras. A otimização dos consequentes por gradiente descendente reduziu significativamente o erro de aproximação~\cite{haykin2008neural}.

A comparação entre diferentes funções de pertinência (gaussiana, triangular, trapezoidal) e operadores (produto, mínimo) mostrou que a escolha da função de pertinência tem impacto relevante no desempenho, enquanto a diferença entre operadores foi pequena neste caso~\cite{ross2010fuzzy}.

Aumentar o número de regras, ajustar os parâmetros das funções de pertinência e otimizar os consequentes são estratégias eficazes para melhorar a aproximação. O sistema fuzzy Takagi-Sugeno mostrou-se uma ferramenta poderosa para modelagem de funções não lineares.

Ambos os sistemas podem ser aprimorados com a inclusão de mais variáveis, regras ou técnicas de otimização. Para aplicações reais, recomenda-se validar o modelo com dados experimentais e considerar fatores externos não modelados.